\section*{Chapitre \ref{ch:ll:linux-tuning} --- Réglage du noyau Linux}

\begin{solution}{exo:ll:linux-tuning:1}
$1\,000 \times \qty{2}{\micro\second} = 0.2~\%$ du temps. L'attente moyenne ajoutée vaut $1\,000 \times (\qty{2}{\micro\second})^2 /
(2 \times \qty{1}{\second}) = \qty{2}{\nano\second}$. Sur un processeur sans tick qui n'exécute qu'une tâche, le tick n'interrompt plus
le thread~: le pire cas perd la durée complète du tick, ici \qty{2}{\micro\second}.
\end{solution}

\begin{solution}{exo:ll:linux-tuning:2}
L'appel échoue avec \texttt{ENOMEM}~: sur ce portable, la limite de mémoire verrouillable d'un utilisateur ordinaire est de 64~MiB,
et 3~GiB la dépasse largement. Il faut relever \texttt{RLIMIT\_MEMLOCK} pour le processus (dans la configuration du service ou
dans les limites du système), ou lui accorder la capacité \texttt{CAP\_IPC\_LOCK}, pour laquelle la limite n'est pas appliquée~; la règle \texttt{memlock} de l'audit vérifie la première.
\end{solution}

\begin{solution}{exo:ll:linux-tuning:3}
\texttt{isolcpus=12-15,44-47} et \texttt{nohz\_full=12-15} (avec \texttt{rcu\_nocbs=12-15}). Les jumeaux 44 à 47 partagent les
unités d'exécution des cœurs chauds~: il faut soit les mettre hors ligne, soit les laisser isolés sans rien qui y soit épinglé,
afin qu'aucun autre thread ne rivalise avec les threads chauds pour leurs cœurs.
\end{solution}

\begin{solution}{exo:ll:linux-tuning:4}
Les deux volent 1~\%. Les mille courtes pauses ajoutent une attente moyenne de $1\,000 \times (10^{-5})^2/2 = \qty{50}{\nano\second}$, et 0.5~\%
des messages attendent plus de \qty{5}{\micro\second}. La pause unique de \qty{10}{\milli\second} ajoute $(10^{-2})^2/2$ secondes,
soit \qty{50}{\micro\second} en moyenne, mille fois plus, et près de 1~\% des messages attendent plus de
\qty{5}{\micro\second}, certains pendant dix millisecondes.
\end{solution}

\begin{solution}{exo:ll:linux-tuning:5}
Lorsque les microsecondes gagnées valent plus qu'un cœur, on paie sa puissance et sa chaleur~: sur les quelques threads du chemin critique, et non sur le reste.
Outre le cœur, cela coûte de la puissance et de la chaleur (la documentation du noyau signale que l'interrogation dans la boucle d'inactivité peut surchauffer le processeur et désactive les fréquences turbo),
et sur un cœur à hyper-threading cela ralentit ce qui tourne sur le jumeau.
\end{solution}

\begin{solution}{exo:ll:linux-tuning:6}
Avec le throttling par défaut, les threads temps réel peuvent utiliser \qty{950}{\milli\second} de chaque seconde~; lorsque l'agent est prêt à s'exécuter, il obtient
les \qty{50}{\milli\second} restantes~: une pause allant jusqu'à \qty{50}{\milli\second} une fois par seconde, soit 5~\% du temps. Corrections~: sortir l'agent du processeur (l'isoler),
après quoi le throttling n'a plus rien à céder~; ou désactiver le throttling, au risque de bloquer le processeur
si le thread se comporte mal. Mieux vaut encore exécuter le thread du flux comme un thread ordinaire sur un cœur isolé.
\end{solution}

\begin{solution}{exo:ll:linux-tuning:7}
Analyser dans la fixture \texttt{/proc/interrupts} les lignes dont le nom contient l'interface (par exemple \texttt{ens1f0-TxRx-0}),
et exiger que le \texttt{smp\_affinity\_list} de chacune soit un sous-ensemble des processeurs de service du nœud du plan. La
fixture réglée les lie aux processeurs 16 à 20 et réussit~; ajouter à la fixture mal réglée une ligne liée au processeur 3 (nœud 0) et
vérifier que l'échec la nomme.
\end{solution}

\begin{solution}{exo:ll:linux-tuning:8}
Les processeurs isolés sont retirés de l'équilibrage de charge de l'ordonnanceur~: les threads dont l'affinité couvre plusieurs processeurs isolés ne sont pas
répartis entre eux et peuvent tous rester sur celui où ils ont démarré, en se partageant ce seul processeur (la courbe commune de la
\cref{fig:ll:linux-tuning:hiccups}). Il faut épingler chaque thread à son propre processeur isolé, comme le fait \texttt{firm.affinity} à partir du plan.
\end{solution}

\begin{solution}{pb:ll:linux-tuning:1}
\begin{enumerate}
\item Les deux amèneraient le noyau sur le processeur (un \omterm{def:ll:cpp-for-latency-i-memory:fault}{défaut de page}, un appel système), et le mesureur se mesurerait
lui-même au lieu de mesurer ce qui lui est imposé.
\item Sur la mesure retenue, 5 secondes pour environ 323 millions de passages~: environ \qty{15}{\nano\second} par passage, donc
une pause de \qty{200}{\nano\second} représente plus de dix passages perdus, et non une itération lente.
\item Environ 1.3~\%, avec une pause maximale d'environ \qty{5.4}{\milli\second}.
\item Environ 51~\%~: deux threads identiques qui demandent toujours le processeur en obtiennent chacun la moitié.
\item Environ \qty{3.1}{\micro\second} sur la mesure retenue~; cette valeur varie avec la pause la plus longue de l'exécution.
\item Environ 0.45~\% au-delà de \qty{10}{\micro\second}, et environ 0.15~\% au-delà de \qty{100}{\micro\second}.
\item Entre \qty{10}{\micro\second} et \qty{100}{\micro\second}~: 50.2~\% des messages attendent plus que le premier seuil et
49.0~\% plus que le second. Environ la moitié n'attend rien, et ceux qui attendent attendent des millisecondes.
\item L'attente croît avec le carré d'une pause~: un message a plus de chances de tomber dans une longue pause, et il attend alors plus longtemps.
\item $250 \times \qty{2}{\micro\second} = 0.05~\%$ et $250 \times (\qty{2}{\micro\second})^2 / 2 = \qty{0.5}{\nano\second}$.
\item $250 \times (\qty{2}{\micro\second} - \qty{1}{\micro\second}) / \qty{1}{\second} = 0.025~\%$.
\item $2 \times 10^{-6}$ du temps et \qty{0.002}{\nano\second} en moyenne.
\item Sur un processeur de service, qui l'exécute pour le compte des processeurs sans tick.
\item Jusqu'à \qty{50}{\milli\second} une fois par seconde, soit 5~\% du temps~: pire que le portable non réglé.
\item Il s'agit d'une machine virtuelle~: son noyau n'a pas de ticks adaptatifs, le travail propre de l'hyperviseur vole du temps
que l'invité ne voit pas ni ne peut déplacer, et un utilisateur ordinaire ne peut pas modifier la ligne de démarrage.
\item \textbf{Résultat retenu.} Épinglé seul sur ce portable, un thread en \omterm{def:ll:linux-tuning:polling}{attente active} ne s'exécute pas environ 1.3~\% du temps, avec une
pause maximale d'environ \qty{5.4}{\milli\second}~; sur l'hôte réglé modélisé, $2 \times 10^{-6}$ du temps, avec une pause maximale de
\qty{2}{\micro\second}~: environ six mille fois moins de temps perdu, et une pause maximale environ 2\,700 fois plus courte.
\item \texttt{isolated}, \texttt{siblings}, \texttt{nohz\_full}, \texttt{rcu\_nocbs}, \texttt{irq\_affinity},
\texttt{cstate} et \texttt{memlock}~: sept sur onze.
\item Sur un processeur isolé de réserve du même hôte, configuré comme les processeurs chauds, ou sur un processeur chaud avant
l'ouverture~; jamais à côté d'un thread chaud.
\item Sa cause~: le mesureur voit que le processeur a été pris, sans savoir par qui. La trace du noyau (compteurs d'interruptions
par processeur dans \texttt{/proc/interrupts} avant et après, événements de trace de l'ordonnanceur) permet de l'attribuer.
\item Le rapport d'audit de l'hôte, à côté du plan de placement, la version du noyau et sa ligne de démarrage.
\item Il donne au chemin critique des processeurs entiers que le noyau laisse tranquilles, au prix des cœurs, de la puissance et
d'un hôte à configurer et à contrôler pour cela.
\end{enumerate}
\end{solution}

\begin{solution}{iq:ll:linux-tuning:1}
\texttt{isolcpus} retire des processeurs à l'équilibrage de charge, de sorte que seuls les threads épinglés explicitement y tournent~; \texttt{nohz\_full} arrête
le tick de l'horloge d'ordonnancement sur ces processeurs lorsqu'ils n'exécutent qu'une tâche~; \texttt{rcu\_nocbs} déplace hors de ces processeurs les rappels différés du noyau.
Chacun supprime une source d'interruption différente, et le tick ne peut s'arrêter tant que d'autres tâches ou rappels ont encore besoin du processeur.
\iqlookfor{trois sources de bruit et la raison pour laquelle chacune demande son propre réglage.}
\end{solution}

\begin{solution}{iq:ll:linux-tuning:2}
Une périodicité régulière évoque le tick ou un minuteur~: vérifier la fréquence de tick du noyau et si le processeur est
sans tick~; comparer les compteurs d'interruptions par processeur avant et après une exécution~; chercher d'autres tâches prêtes sur ce processeur et le travail du noyau par processeur~; tracer
les événements de l'ordonnanceur sur ce processeur. Confirmer avec un mesureur de pauses sur le même processeur.
\iqlookfor{mesurer d'abord, puis attribuer la cause avec des compteurs et du traçage.}
\end{solution}

\begin{solution}{iq:ll:linux-tuning:3}
Le blocage ne coûte rien à l'inactivité, mais ajoute le réveil, des dizaines de microsecondes sur la machine du chapitre. L'\omterm{def:ll:linux-tuning:polling}{attente active} coûte
un cœur qui tourne à vide et économise le réveil (\qty{2}{\micro\second} contre \qty{30}{\micro\second} ici). Le \omterm{def:ll:linux-tuning:polling}{contournement du noyau} coûte
une carte compatible, une pile réseau séparée et un cœur qui tourne à vide, et retire le noyau du chemin. Le choix se fait selon le \omterm{def:ll:where-latency-comes-from:budget}{budget de latence},
thread par thread.
\iqlookfor{les coûts de chaque côté et le budget comme règle de décision.}
\end{solution}

\begin{solution}{iq:ll:linux-tuning:4}
En général non, s'ils tournent seuls sur des cœurs isolés~: il n'y a alors rien à préempter. Les risques sont réels~: le throttling temps réel (une pause de
\qty{50}{\milli\second} par seconde par défaut lorsqu'une autre tâche est prête), un thread emballé qui bloque un processeur, et l'\omterm{def:ll:lock-free-programming:inversion}{inversion de priorité} avec des threads de priorité plus basse avec lesquels il partage des verrous
(chapitre~11).
\iqlookfor{throttling et famine, et non seulement «~plus de priorité, donc plus de rapidité~».}
\end{solution}

\begin{solution}{iq:ll:linux-tuning:5}
Sur les processeurs de service du socket de la carte~: hors des cœurs chauds, pour que ceux-ci ne soient pas interrompus, mais sur le nœud de
la carte, afin que le travail du noyau sur le paquet reste dans le cache partagé du socket, près des threads qui le lisent, et ne
traverse pas l'interconnexion.
\iqlookfor{isolement par rapport au cœur chaud, proximité de la carte.}
\end{solution}

\begin{solution}{iq:ll:linux-tuning:6}
Les rares et longues pauses que le réglage n'a pas supprimées~: des interruptions encore autorisées sur les processeurs chauds, un tick qui tourne encore (une
deuxième tâche prête à s'exécuter, un \omterm{def:ll:linux-tuning:nohz}{noyau sans ticks} adaptatifs), des \omterm{def:ll:cpp-for-latency-i-memory:fault}{défauts de page} (mémoire non verrouillée ou non préchargée), un jumeau d'hyper-threading utilisé,
l'hyperviseur sur une machine virtuelle. Mesurer les pauses sur le processeur chaud et auditer l'hôte.
\iqlookfor{la queue vient d'événements rares~; les carrés de la proposition.}
\end{solution}
